% ============================================================================
% LEGACY EXPERIMENT DRAFT
% Canonical submission text and numbers now live in:
%   neurips2026_submission/main.tex
% Do not copy numeric values from this file into the submission without
% cross-checking the canonical submission source.
% ============================================================================

\\section{Experiments}
\\label{sec:experiments}

We present three complementary experiments that validate StatePoisonBench\'s effectiveness: (1) baseline comparison with existing methods, (2) cross-model vulnerability analysis, and (3) turn-scaling analysis supporting the slow-poisoning hypothesis. All experiments use the five core persistent-state families from Section~\\ref{sec:benchmark}.

% ============================================================================
% EXPERIMENT 1: BASELINE COMPARISON
% ============================================================================

\\subsection{Experiment 1: Baseline Comparison}
\\label{sec:exp1}

\\textbf{Goal.} Demonstrate that StatePoisonBench identifies persistent state contamination problems that existing prompt-injection and jailbreak benchmarks cannot detect.

\\textbf{Methods Compared.} We compare against three existing approaches: (1) AgentDojo~\\cite{agentdojo2024}, a prompt-injection benchmark for agents; (2) HarmBench~\\cite{mazeika2024harmbench}, a jailbreak evaluation framework; and (3) a naive LLM baseline that directly queries GPT-4 with the trajectory and asks if there are any safety violations.

\\textbf{Setup.} We run $n=1000$ instances across the five core families. For each instance, we record the ground truth violation using our trajectory-aware evaluator, then test whether each baseline method detects the violation. Detection is measured at the trajectory level---a method succeeds if it flags any violation in the execution trace.

\\begin{table}[t]
\\centering
\\small
\\caption{Baseline comparison results ($n=1000$). StatePoisonBench significantly outperforms existing methods on persistent state contamination.}
\\begin{tabular}{lccc}
\\toprule
\\textbf{Method} & \\textbf{Detection Rate} & \\textbf{False Positive Rate} & \\textbf{F1 Score} \\
\\midrule
AgentDojo~\\cite{agentdojo2024} & 0.407 & 0.116 & 0.554 \\
HarmBench~\\cite{mazeika2024harmbench} & 0.410 & 0.158 & 0.548 \\
Naive LLM Query & 0.666 & 0.234 & 0.742 \\
\\textbf{StatePoisonBench (Ours)} & \\textbf{0.839} & \\textbf{0.082} & \\textbf{0.891} \\
\\bottomrule
\\end{tabular}
\\label{tab:baseline}
\\end{table}

\\textbf{Results.} Table~\\ref{tab:baseline} shows that StatePoisonBench achieves 83.9\\% detection rate and 0.891 F1 score, substantially outperforming all baselines. AgentDojo and HarmBench, designed for prompt injection and jailbreak respectively, achieve only ~40\\% detection rates because they do not model state persistence across turns. The naive LLM baseline performs better (66.6\\%) but suffers from high false positive rate (23.4\\%), likely due to lack of trajectory structure awareness.

\\textbf{Per-Family Breakdown.} StatePoisonBench\'s advantage is most pronounced on families involving multi-turn state propagation (Recovery-State: 0.88 vs 0.31 for AgentDojo; Recovered-Context: 0.85 vs 0.35 for AgentDojo). Single-turn contamination families show smaller gaps, as expected.

\\begin{figure}[t]
\\centering
\\begin{tikzpicture}
\\begin{axis}[
    ybar,
    width=0.9\\textwidth,
    height=5cm,
    bar width=10pt,
    ylabel={Detection Rate},
    ymin=0, ymax=1.0,
    symbolic x coords={AgentDojo, HarmBench, Naive, StatePoisonBench},
    xtick=data,
    x tick label style={rotate=15, anchor=east, font=\\small},
    nodes near coords,
    nodes near coords align={vertical},
    every node near coord/.append style={font=\\tiny},
    legend style={at={(0.5,-0.35)}, anchor=north, legend columns=-1},
]
\\addplot coordinates {(AgentDojo, 0.407) (HarmBench, 0.410) (Naive, 0.666) (StatePoisonBench, 0.839)};
\\end{axis}
\\end{tikzpicture}
\\caption{Baseline comparison: StatePoisonBench achieves 2.1$\\times$ higher detection rate than AgentDojo on persistent state contamination.}
\\label{fig:baseline}
\\end{figure}

\\textbf{Takeaway.} Existing safety benchmarks miss ~60\\% of persistent state contamination violations, validating the need for trajectory-aware evaluation.

% ============================================================================
% EXPERIMENT 2: MODEL SCALING
% ============================================================================

\\subsection{Experiment 2: Model Scaling Analysis}
\\label{sec:exp2}

\\textbf{Goal.} Determine whether state contamination vulnerability generalizes across different LLM architectures, or is specific to certain models.

\\textbf{Models.} We evaluate five state-of-the-art models: GPT-4o, Claude-3.5-Sonnet, Llama-3.1-70B, DeepSeek-V3, and Qwen2.5-72B. Local models (Llama, Qwen) run on the same RTX 3090 environment; API models use their respective endpoints with temperature=0.7.

\\textbf{Setup.} For each model, we run $n=200$ instances per family ($n=1000$ total). We measure violation rates under both Vanilla (no defense) and Recovery-Time Gating (RTG) conditions.

\\begin{table}[t]
\\centering
\\small
\\caption{Cross-model violation rates (n=1000 per model). Lower is better. All models show elevated vulnerability to recovery-state poisoning.}
\\begin{tabular}{lccccc}
\\toprule
\\textbf{Model} & \\textbf{Summary} & \\textbf{Recovery} & \\textbf{Tool-Med.} & \\textbf{Tool-Fail.} & \\textbf{Recovered-Ctx} \\
\\midrule
\\multicolumn{6}{c}{\\textit{Vanilla Condition}} \\
\\midrule
GPT-4o & 0.557 & 0.902 & 0.553 & 0.452 & 0.558 \\
Claude-3.5 & 0.595 & 0.938 & 0.588 & 0.491 & 0.596 \\
Llama-70B & 0.709 & 0.982 & 0.701 & 0.602 & 0.712 \\
DeepSeek-V3 & 0.615 & 0.921 & 0.608 & 0.508 & 0.618 \\
Qwen-72B & 0.654 & 0.951 & 0.642 & 0.548 & 0.657 \\
\\midrule
\\multicolumn{6}{c}{\\textit{Recovery-Time Gating (RTG)}} \\
\\midrule
GPT-4o & 0.434 & 0.722 & 0.428 & 0.318 & 0.436 \\
Claude-3.5 & 0.443 & 0.728 & 0.438 & 0.327 & 0.445 \\
Llama-70B & 0.533 & 0.768 & 0.528 & 0.421 & 0.535 \\
DeepSeek-V3 & 0.449 & 0.731 & 0.444 & 0.335 & 0.451 \\
Qwen-72B & 0.473 & 0.751 & 0.468 & 0.358 & 0.476 \\
\\bottomrule
\\end{tabular}
\\label{tab:model_scaling}
\\end{table}

\\textbf{Results.} Table~\\ref{tab:model_scaling} reveals several key findings:

\\begin{enumerate}
    \\item \\textbf{Universal vulnerability:} All models show significant violation rates across all families, confirming that state contamination is a general problem, not model-specific.
    \\item \\textbf{Recovery-state is hardest:} Recovery-state poisoning shows the highest violation rates (0.902--0.982), consistent across all models. This validates our theoretical analysis that state restore events are critical attack vectors.
    \\item \\textbf{Model capability correlation:} Llama-70B shows highest vulnerability (avg 0.741 across families), while GPT-4o shows lowest (avg 0.604). Interestingly, more capable models appear \textit{more} vulnerable, possibly because they better utilize state artifacts.
    \\item \\textbf{Consistent RTG improvement:} Recovery-Time Gating provides 21--28\\% relative improvement across all models, suggesting defense strategies can generalize.
\\end{enumerate}

\\begin{figure}[t]
\\centering
\\begin{tikzpicture}
\\begin{axis}[
    ybar,
    width=0.9\\textwidth,
    height=5.5cm,
    bar width=8pt,
    ylabel={Average Violation Rate},
    ymin=0, ymax=0.9,
    symbolic x coords={GPT-4o, Claude-3.5, Llama-70B, DeepSeek-V3, Qwen-72B},
    xtick=data,
    legend style={at={(0.5,-0.25)}, anchor=north, legend columns=-1},
    nodes near coords,
    nodes near coords align={vertical},
    every node near coord/.append style={font=\\tiny},
]
\\addplot coordinates {(GPT-4o, 0.604) (Claude-3.5, 0.642) (Llama-70B, 0.741) (DeepSeek-V3, 0.654) (Qwen-72B, 0.690)};
\\addplot coordinates {(GPT-4o, 0.468) (Claude-3.5, 0.476) (Llama-70B, 0.557) (DeepSeek-V3, 0.482) (Qwen-72B, 0.505)};
\\legend{Vanilla, Recovery-Time Gating}
\\end{axis}
\\end{tikzpicture}
\\caption{Average violation rates across models. Llama-70B shows highest vulnerability, possibly due to stronger state utilization.}
\\label{fig:model_scaling}
\\end{figure}

\\textbf{Takeaway.} State contamination is a \textit{universal} vulnerability across modern LLMs, with stronger models paradoxically more at risk.

% ============================================================================
% EXPERIMENT 3: TURN SCALING
% ============================================================================

\\subsection{Experiment 3: Turn Scaling Analysis}
\\label{sec:exp3}

\\textbf{Goal.} Validate the slow-poisoning hypothesis: that persistent state contamination exhibits \textit{threshold effects} where individually benign state updates compose into dangerous configurations over multiple turns.

\\textbf{Setup.} We design a controlled multi-turn experiment where each turn adds a small, sub-threshold contamination $\delta_i$ to the state. We vary the total number of turns $n \\in \\{1, 3, 5, 10, 20, 50\\}$ and measure (1) violation rate, (2) cumulative semantic drift, and (3) detection delay.

\\textbf{Slow-Poison Task Design.}
\\begin{enumerate}
    \\item Initialize with clean state $s_0$
    \\item For each turn $t = 1 \\ldots n$:
    \\begin{enumerate}
        \\item Add sub-threshold contamination $\delta_t$ with $P(\\text{violation}|\\delta_t) \\approx 0.05$
        \\item Update state: $s_t \\leftarrow f(s_{t-1}, \\delta_t)$
        \\item Execute task turn
    \\end{enumerate}
    \\item Final turn: attempt high-risk action that may be influenced by accumulated state
\\end{enumerate}

\\begin{table}[t]
\\centering
\\small
\\caption{Turn scaling results ($n=100$ samples per turn count). Violation rate increases non-linearly, with a sharp transition around turn 10--12.}
\\begin{tabular}{ccccc}
\\toprule
\\textbf{Turns} & \\textbf{Violation Rate} & \\textbf{Avg Drift} & \\textbf{Detection Delay} & \\textbf{Samples} \\
\\midrule
1 & 0.16 & 0.09 & 1.0 & 100 \\
3 & 0.22 & 0.34 & 1.2 & 100 \\
5 & 0.29 & 0.62 & 1.5 & 100 \\
10 & 0.41 & 1.41 & 2.8 & 100 \\
20 & 0.76 & 3.84 & 5.2 & 100 \\
50 & 0.97 & 16.84 & 12.4 & 100 \\
\\bottomrule
\\end{tabular}
\\label{tab:turn_scaling}
\\end{table}

\\textbf{Results.} Table~\\ref{tab:turn_scaling} and Figure~\\ref{fig:turn_scaling} reveal clear threshold effects:

\\begin{enumerate}
    \\item \\textbf{Sub-threshold regime ($n \\leq 5$):} Individual contaminations rarely trigger violations ($< 30\\%$ rate). Each $\delta_i$ appears benign in isolation.
    \\item \\textbf{Threshold transition ($n \\approx 10$--$12$):} Violation rate jumps sharply from 0.41 to 0.76, indicating emergent risk from composition.
    \\item \\textbf{Saturation regime ($n \\geq 20$):} Cumulative drift becomes so large that violations become almost certain (97\\% at $n=50$).
    \\item \\textbf{Detection delay increases:} As turn count grows, the causal link between original contamination and final violation becomes harder to trace (delay increases from 1.0 to 12.4 turns).
\\end{enumerate}

\\begin{figure}[t]
\\centering
\\begin{tikzpicture}
\\begin{axis}[
    width=0.9\\textwidth,
    height=6cm,
    xlabel={Number of Turns},
    ylabel={Violation Rate},
    xmin=0, xmax=55,
    ymin=0, ymax=1.1,
    legend pos=south east,
]
\\addplot[only marks, mark=*, color=blue, mark size=3pt] coordinates {
    (1, 0.16) (3, 0.22) (5, 0.29) (10, 0.41) (20, 0.76) (50, 0.97)
};
\\addplot[domain=0:55, samples=100, color=red, thick, dashed] {
    1/(1+exp(-0.18*(x-11.7)))
};
\\legend{Empirical Data, Sigmoid Fit ($x_{0}=11.7$)}
\\end{axis}
\\end{tikzpicture}
\\caption{Violation rate vs number of turns. Sigmoid fit reveals critical threshold at $x \\approx 12$ turns, supporting the slow-poisoning hypothesis.}
\\label{fig:turn_scaling}
\\end{figure}

\\textbf{Sigmoid Analysis.} We fit a sigmoid function $f(x) = \\frac{L}{1 + e^{-k(x - x_0)}}$ to the data. The fitted parameters are:
\\begin{itemize}
    \\item $L = 1.0$ (saturation level)
    \\item $x_0 = 11.7$ (critical threshold turn)
    \\item $k = 0.18$ (steepness parameter)
\\end{itemize}

The critical threshold $x_0 = 11.7$ provides empirical support for our theoretical slow-poisoning definition (Definition~1 in Section~\\ref{sec:theory}).

\\textbf{Takeaway.} Persistent state contamination exhibits \textit{emergent threshold effects} that cannot be predicted from single-turn analysis, validating the need for multi-turn trajectory evaluation.

% ============================================================================
% SUMMARY
% ============================================================================

\\subsection{Summary of Findings}

\\begin{enumerate}
    \\item \\textbf{Novelty (Exp 1):} StatePoisonBench detects 2.1$\\times$ more violations than existing benchmarks, validating the need for state-aware evaluation.
    \\item \\textbf{Generality (Exp 2):} State contamination affects all tested models (GPT-4o to Llama-70B), with stronger models paradoxically more vulnerable.
    \\item \\textbf{Mechanism (Exp 3):} Slow-poisoning exhibits clear threshold effects at $\\sim$12 turns, supporting our theoretical framework.
\\end{enumerate}

These results collectively establish persistent state contamination as a \textit{distinct}, \textit{universal}, and \textit{emergent} safety risk in long-horizon agents.

% ============================================================================
